\usepackage[bookmarks=true, bookmarksnumbered=true, bookmarksopen=true]{hyperref} % To control the TOC
\usepackage{polyglossia} % To use hebrew
\usepackage{fontspec} % To control hebrew fonts
\usepackage{setspace} % To control line spacing
% Hebrew font chosen: Arial
\newfontfamily\hebrewfont{Arial}[Script=Hebrew]
% Set hebrew line spacing manually
\makeatletter
\g@addto@macro\hebrewfont{\setstretch{1.3}}
\makeatother
% Define hebrew page numbering layout
% Adapted from a Stack Overflow answer by user 'John Kormylo'
% Source: https://tex.stackexchange.com/a/303111/442136
% Licensed under CC-BY-SA 3.0
\makeatletter
\def\ps@hebrew{%
	\let\@oddhead\relax
	\let\@evenhead\relax
	\def\@oddfoot{\hfill\begin{otherlanguage}{hebrew}%
			\hebrewnumeral{\arabic{page}}%
		\end{otherlanguage}\hfill}%
	\let\@evenfoot\@oddfoot}
\makeatother

% Add the hebrew numbered section to TOC
% Adapted from a Stack Overflow answer by user 'user2679290'
% Source: https://tex.stackexchange.com/a/617122/442136
% Licensed under CC-BY-SA 4.0
% It does three things: 
% 1. Replace \thepage with the letter aleph (we need robustify as this is not run directly)
% 2. Adds a section to the TOC with this aleph as page number (and correct link as we use the hyperref version of \addcontentsline)
% 3. Restore \thepage to its last value (we assume the location we called the command is just before the abstract and have alph pagenumbering)
\newcommand{\addhebrewtoctitle}[1]{\robustify{\thepage}
\addtocontents{toc}{\protect\renewcommand\thepage{\texthebrew{\protect\hebrewnumeral{1}}}}
\addcontentsline{toc}{section}{#1}
\addtocontents{toc}{\protect\renewcommand\thepage{\alph{page}}}}